% ============================================================
%  序（frontmatter/preface.tex）
%  由 main.tex \input 引入（见 main.tex 中关于 \includeonly 的注）。
% ============================================================

\chapter*{序　Preface}

% —— 入目录（不编号）——
% 标准 book 的 \@dottedtocline 会给 chapter 级条目留 \numberline 位（1.5em），
% 无编号时该位空着，目录里「序」会比「第一章」左端缩进偏多、看着不平齐。
% 此处临时令 \numberline 无效，使无编号条目与编号条目左端对齐。
\begingroup
    \renewcommand{\numberline}[1]{}   % 局部重定义，出组自动复原
    \addcontentsline{toc}{part}{序　Preface}
\endgroup

% —— 页眉/页脚 ——
% \chapter* 不自动设 \leftmark，故文末显式补上（与普通章的页眉同制）。
\thispagestyle{plain}   % 序首面不置页眉，仅页码，清整

近者同学多以化学基础之题见询。往来问难，日不暇给，而环顾同侪，能裒然成帙，勒成一编者，盖无二三。或曰：化学之途，门径既多，名目尤夥；元素之性，反应之则，计算之方，实验之规，皆初学者所苦。苟无统纪，则如入宝山而迷途，临沧海而失楫。余闻而然之，心自戚戚；思有以助之，而未得其便。是乃发箧中故纸，堆案盈尺，纷然杂陈，爬梳别择，去重补阙，正讹芟杂，存其精要，裒为一编。其有疑义，则参诸旧记；其有阙略，则询诸同窗。虽粗具梗概，未敢自谓成书，然于初学或不无涓埃之补。挂漏之处，尚祈同窗赐以匡正。

是编之例，略以类从：宁简毋滥，宁确毋夸；不敢求详于高深，惟期有裨于初学。

是编所辑，实本　\uline{李}先生\uline{之春}、　\uline{龚}先生\uline{奇峰}、　\uline{扶}先生\uline{云松}之教。诸先生既授其学，复示其方；举凡肯綮所在，多由口授而心得。微言奥旨，既得之于耳；条分缕析，复录之于纸。或剖析于课堂，或指示于言谈；一字之训，必求其安；一义之疑，必究其归。我曹所能，惟写录编次耳。昔人云「述而不作，信而好古」，则此编之纂惟诸先生之述是集，而不敢自谓作。若有毫发之明，则诸先生所赐；其有纰缪，则我曹写录之疏也。

至于图绘，皆出同人手摹，盖文字所不能尽者，图以形之；言语所不能详者，绘以明之。同人各出己意，手自摹写，或临旧本，或依口授，或参己见，务求有助参览，有裨解悟，有参补益。然笔力有工拙，见闻有浅深，非敢自诩精工。点画之间，疏略难免，或标识未周，或笔意粗略，或描摹失真。惟冀览者不以为陋，原其拙而摘其误，或加匡正，或赐教言，则受益多矣，为幸厚矣。

书既竟，欲题数语于卷端，而忽忘所拟之辞。既而思之，夫化者，格物之枢也。天地之间，自金石之顽，至于血气之灵，莫非质之所聚、变之所行。其聚也，有所以聚之理；其散也，有所以散之故。观其聚散之所以然，则化之学理可致也。故其为学，上穷元素之微，下究万物之变；分之则原子离子，合之则反应原理；或化合而新质生，或分解而旧体易；或氧化还原而电子相易，或沉淀溶解而平衡相持。其间数目有计量之严，实验有操作之则，图式有符号之约，皆不可忽也。虽然，学化学者，贵在自得于心，反复推究，实验以验其真，思维以通其变；此编不过一导路之灯耳。若徒诵记而不思，则如观棋谱而不弈；若畏难而中辍，则如望山巅而不登。愿诸君学之日进，如催化之于\(E_\text{a}\)，行之愈速；如滴定之近突跃，积之愈谨；如结晶之渐析出，得之愈纯。更上一层，以尽其妙。

是为序。

% —— 题识（小字，右齐）与落款（再收半级）——
% 题识为「识（zhì）」体：不尽以书名归诸先生，只自任写录之劳，
% 版面上令读者最后看到的是三位先生的姓名。
\vspace{1.5\baselineskip}
{\small\raggedleft
是编文字，多本　\uline{李}先生\uline{之春}、　\uline{龚}先生\uline{奇峰}、　\uline{扶}先生\uline{云松}之教。\par
既录其言，不敢没其赐，因书于端，以识所自。\par
}
\vspace{1.2\baselineskip}
{\footnotesize\raggedleft
本书编写组　谨识\par
}

\markboth{序　Preface}{}
